% GENERATED FILE. Edit canonical lessons, then run npm run generate:books.
% canonical-source-hash: fnv1a64:ec30d3288af49a09
% canonical-lessons: FR-C118-bijou, FR-C118-bague, FR-C118-collier, FR-C118-parfum, FR-C118-lettre

\chapter{Jewel, Ring, Necklace, Perfume, Letter}
\label{ch:fr-jewel-ring-necklace-perfume-letter}
\hlchaptermodality{118}

\begin{chapteropening}
\textbf{By the end of this chapter:} \emph{I can name a jewel, a ring, a necklace, perfume, a letter.}

Complete the last lesson of chapter 118: I can name a jewel, a ring, a necklace, perfume, a letter.
\end{chapteropening}

\section[le bijou]{le bijou — a jewel, a piece of jewellery}

\label{lesson:FR-C118-bijou}



\begin{quote}
Before the new one: say the French for size; the waist, then the French for a wallet.
\end{quote}

\subsection*{You'll want to know: le bijou}
\textbf{le bijou} — \textquotedblleft{}a jewel, a piece of jewellery\textquotedblright{}. Plural \textbf{les bijoux}, with an x.

\begin{cousinweb}
\textbf{bijou} $\leftarrow$ Breton \emph{bizou}, \textquotedblleft{}a ring\textquotedblright{}, from \emph{biz}, \textquotedblleft{}finger\textquotedblright{} — borrowed. (EtymDB 2.1)
\end{cousinweb}

\begin{grammarlens}[title={What you've built}]
One more word for this chapter.
\end{grammarlens}

\subsection*{Your turn}
\begin{itemize}
\raggedright
  \item \emph{Say it:} \emph{le bijou}
  \item \emph{Say it:} \emph{le bijou}, once more
  \item \emph{Say it:} say \emph{le bijou}
  \item \emph{Recall:} say the French for size; the waist, then the French for a wallet, then say \emph{le bijou} again
\end{itemize}

\begin{tcolorbox}[breakable,colback=teal!4,colframe=teal!35!black,title={Before you move on}]
What does le bijou mean? (\textquotedblleft{}A jewel, a piece of jewellery\textquotedblright{}.) Say it once more.
\end{tcolorbox}
\section[la bague]{la bague — a ring}

\label{lesson:FR-C118-bague}



\begin{quote}
Before the new one: say the French for a wallet, then the French for a jewel.
\end{quote}

\subsection*{You'll want to know: la bague}
\textbf{la bague} — \textquotedblleft{}a ring\textquotedblright{}.

\begin{cousinweb}
\textbf{bague} — its origin is uncertain; Middle Dutch \emph{bagge}, \textquotedblleft{}a ring\textquotedblright{}, is the usual guess. (EtymDB 2.1 has its Middle French form.)
\end{cousinweb}

\begin{grammarlens}[title={What you've built}]
One more word for this chapter.
\end{grammarlens}

\subsection*{Your turn}
\begin{itemize}
\raggedright
  \item \emph{Say it:} \emph{la bague}
  \item \emph{Say it:} \emph{la bague}, once more
  \item \emph{Say it:} say \emph{la bague}
  \item \emph{Recall:} say the French for a wallet, then the French for a jewel, then say \emph{la bague} again
\end{itemize}

\begin{tcolorbox}[breakable,colback=teal!4,colframe=teal!35!black,title={Before you move on}]
What does la bague mean? (\textquotedblleft{}A ring\textquotedblright{}.) Say it once more.
\end{tcolorbox}
\section[le collier]{le collier — a necklace}

\label{lesson:FR-C118-collier}



\begin{quote}
Before the new one: say the French for a jewel, then the French for a ring.
\end{quote}

\subsection*{You'll want to know: le collier}
\textbf{le collier} — \textquotedblleft{}a necklace\textquotedblright{}.

\begin{cousinweb}
\textbf{collier} $\leftarrow$ Late Latin \emph{collare}, \textquotedblleft{}a band for the neck\textquotedblright{}, from \emph{collum}, \textquotedblleft{}neck\textquotedblright{} (EtymDB 2.1) — inherited: the same \emph{collum} as \emph{cou}.
\end{cousinweb}

\begin{grammarlens}[title={What you've built}]
One more word for this chapter.
\end{grammarlens}

\subsection*{Your turn}
\begin{itemize}
\raggedright
  \item \emph{Say it:} \emph{le collier}
  \item \emph{Say it:} \emph{le collier}, once more
  \item \emph{Say it:} say \emph{le collier}
  \item \emph{Recall:} say the French for a jewel, then the French for a ring, then say \emph{le collier} again
\end{itemize}

\begin{tcolorbox}[breakable,colback=teal!4,colframe=teal!35!black,title={Before you move on}]
What does le collier mean? (\textquotedblleft{}A necklace\textquotedblright{}.) Say it once more.
\end{tcolorbox}
\section[le parfum]{le parfum — perfume}

\label{lesson:FR-C118-parfum}



\begin{quote}
Before the new one: say the French for a ring, then the French for a necklace.
\end{quote}

\subsection*{You'll want to know: le parfum}
\textbf{le parfum} — \textquotedblleft{}perfume\textquotedblright{}.

\begin{cousinweb}
\textbf{parfum} — formed from the verb for \textquotedblleft{}to scent\textquotedblright{}, which came from Italian \emph{perfumare}, \textquotedblleft{}to smoke through\textquotedblright{}, from Latin \emph{per}, \textquotedblleft{}through\textquotedblright{}, and \emph{fumare}, \textquotedblleft{}to smoke\textquotedblright{}. (EtymDB 2.1 records the Latin parts.)
\end{cousinweb}

\begin{grammarlens}[title={What you've built}]
One more word for this chapter.
\end{grammarlens}

\subsection*{Your turn}
\begin{itemize}
\raggedright
  \item \emph{Say it:} \emph{le parfum}
  \item \emph{Say it:} \emph{le parfum}, once more
  \item \emph{Say it:} say \emph{le parfum}
  \item \emph{Recall:} say the French for a ring, then the French for a necklace, then say \emph{le parfum} again
\end{itemize}

\begin{tcolorbox}[breakable,colback=teal!4,colframe=teal!35!black,title={Before you move on}]
What does le parfum mean? (\textquotedblleft{}Perfume\textquotedblright{}.) Say it once more.
\end{tcolorbox}
\section[la lettre]{la lettre — a letter}

\label{lesson:FR-C118-lettre}



\begin{quote}
Before the new one: say the French for a necklace, then the French for perfume.
\end{quote}

\subsection*{You'll want to know: la lettre}
\textbf{la lettre} — \textquotedblleft{}a letter\textquotedblright{}.

\begin{cousinweb}
\textbf{lettre} $\leftarrow$ Old French \emph{letre} $\leftarrow$ Latin \emph{littera}, \textquotedblleft{}a letter of the alphabet\textquotedblright{} — inherited. (EtymDB 2.1)
\end{cousinweb}

\begin{grammarlens}[title={What you've built}]
One more word for this chapter.
\end{grammarlens}

\subsection*{Your turn}
\begin{itemize}
\raggedright
  \item \emph{Say it:} \emph{la lettre}
  \item \emph{Say it:} \emph{la lettre}, once more
  \item \emph{Say it:} say \emph{la lettre}
  \item \emph{Recall:} say the French for a necklace, then the French for perfume, then say \emph{la lettre} again
\end{itemize}

\begin{tcolorbox}[breakable,colback=teal!4,colframe=teal!35!black,title={Before you move on}]
What does la lettre mean? (\textquotedblleft{}A letter\textquotedblright{}.) Say it once more.
\end{tcolorbox}
